\section{Wstęp}

\subsection{Problem}

Automatyczne wykrywanie fałszywych wiadomości (\emph{fake-news}) to zadanie
klasyfikacji binarnej: na podstawie treści artykułu prasowego model rozstrzyga,
czy pochodzi on z wiarygodnego źródła (\emph{Real}), czy jest dezinformacją
(\emph{Fake}). Trudność problemu nie leży w samym dopasowaniu klasyfikatora ---
jak pokażemy, klasy są w tym zbiorze silnie rozróżnialne leksykalnie --- lecz w
\textbf{rzetelnym przygotowaniu danych}: zbiór zawiera liczne sygnały wycieku
informacji (ang. \emph{data leakage}), które pozwalają uzyskać niemal perfekcyjne
metryki bez uczenia się treści. Większość pracy projektowej dotyczy więc
identyfikacji i usunięcia tych sygnałów, tak by model uczył się stylu i tematyki, a
nie artefaktów technicznych.

\subsection{Cel}

Celem tego etapu jest zbudowanie \textbf{pierwszej sieci neuronowej od zera}
(bez wag pretrenowanych, bez fine-tuningu i feature-tuningu --- zgodnie z wymogiem
przedmiotu SSN) oraz zweryfikowanie kompletnej pętli treningowej w PyTorchu
(propagacja wsteczna, optymalizator Adam, early stopping, regularyzacja, obsługa
GPU). Stanowi to fundament pod model docelowy pracy dyplomowej:
\textbf{BiLSTM + Attention} z osadzeniami GloVe~300d.

\subsection{Zbiór danych}

Wykorzystujemy publiczny zbiór \emph{Fake and Real News Dataset} (Kaggle,
\texttt{clmentbisaillon/fake-and-real-news-dataset}) --- dwa pliki CSV
(\texttt{True.csv}, \texttt{Fake.csv}) z artykułami prasowymi w języku angielskim.
W całym projekcie obowiązuje konwencja etykiet:
\[
  \texttt{1 = Real}, \qquad \texttt{0 = Fake}.
\]
Model operuje wyłącznie na kolumnie \texttt{text}; kolumny \texttt{title},
\texttt{subject} i \texttt{date} są odrzucane jako źródła wycieku informacji
(uzasadnienie w sekcji~\ref{sec:dane}).
